
\documentclass{article}
\usepackage{colortbl}
\usepackage{makecell}
\usepackage{multirow}
\usepackage{supertabular}

\begin{document}

\newcounter{utterance}

\centering \large Interaction Transcript for game `hot\_air\_balloon', experiment `air\_balloon\_survival\_lt\_negotiation\_easy', episode 3 with nemotron{-}3{-}ultra.
\vspace{24pt}

{ \footnotesize  \setcounter{utterance}{1}
\setlength{\tabcolsep}{0pt}
\begin{supertabular}{c@{$\;$}|p{.15\linewidth}@{}p{.15\linewidth}p{.15\linewidth}p{.15\linewidth}p{.15\linewidth}p{.15\linewidth}}
    \# & $\;$A & \multicolumn{4}{c}{Game Master} & $\;\:$B\\
    \hline

    \theutterance \stepcounter{utterance}  
    & & \multicolumn{4}{p{0.6\linewidth}}{
        \cellcolor[rgb]{0.9,0.9,0.9}{
            \makecell[{{p{\linewidth}}}]{
                \texttt{\tiny{[P1$\langle$GM]}}
                \texttt{Dalyvaujate bendradarbiavimo derybų žaidime.} \\
\\ 
\texttt{Kartu su kitu dalyviu turite susitarti dėl vieno daiktų rinkinio, kuris bus išlaikytas. Kiekvienas iš jūsų turi savo nuomonę apie tai, kiek kiekvienas daiktas jums svarbus (svarba). Jūs nežinote, kaip kitas dalyvis vertina daiktus. Be to, jums nurodoma, kiek pastangų reikalauja kiekvienas daiktas.} \\
\texttt{Dėl rinkinio galite susitarti tik tuo atveju, jei pasirinktų daiktų bendros pastangos neviršija bendros ribos:} \\
\\ 
\texttt{LIMIT = 3712} \\
\\ 
\texttt{Čia pateikiamos atskirų daiktų pastangų vertės:} \\
\\ 
\texttt{Daikto pastangos = \{"C08": 949, "C55": 388, "A43": 79, "A67": 585, "B42": 564, "A62": 230, "B19": 580, "C43": 84, "C73": 976, "C10": 274, "B36": 374, "A18": 913, "C61": 303, "C67": 578, "C42": 548\}} \\
\\ 
\texttt{Čia pateikiama jūsų asmeninė nuomonė apie kiekvieno daikto svarbą:} \\
\\ 
\texttt{Daikto svarbos reikšmės = \{"C08": 138, "C55": 583, "A43": 868, "A67": 822, "B42": 783, "A62": 65, "B19": 262, "C43": 121, "C73": 508, "C10": 780, "B36": 461, "A18": 484, "C61": 668, "C67": 389, "C42": 808\}} \\
\\ 
\texttt{Tikslas:} \\
\\ 
\texttt{Jūsų tikslas {-} susitarti dėl bendro daiktų rinkinio, kuris būtų jums kuo naudingesnis (t. y. maksimaliai padidintų bendrą svarbą JUMS), neviršijant LIMIT. Neprivalote kiekvienoje žinutėje pateikti PASIŪLYMAS {-} galite tiesiog derėtis. Visos taktikos yra leidžiamos!} \\
\\ 
\texttt{Sąveikos protokolas:} \\
\\ 
\texttt{Žinutėje galite naudoti tik toliau nurodytus struktūrinius formatus:} \\
\\ 
\texttt{PASIŪLYMAS: \{'A', 'B', 'C', …\}} \\
\texttt{Siūlykite išlaikyti būtent šiuos daiktus.} \\
\texttt{ATSISAKAU: \{'A', 'B', 'C', …\}} \\
\texttt{Aiškiai atmeskite oponento pasiūlymą.} \\
\texttt{ARGUMENTAS: \{'...'\}} \\
\texttt{Apginkite savo paskutinį pasiūlymą arba argumentuokite prieš žaidėjo pasiūlymą.} \\
\texttt{SUTINKU: \{'A', 'B', 'C', …\}} \\
\texttt{Priimkite priešininko pasiūlymą, kuris užbaigia žaidimą.} \\
\\ 
\\ 
\texttt{Taisyklės:} \\
\\ 
\texttt{Galite SUTINKU tik su pasiūlymu, kurį kita šalis užregistravo per PASIŪLYMAS.} \\
\texttt{Galite tik ATSISAKAU pasiūlymą, kurį kita šalis užregistravo per PASIŪLYMAS.} \\
\texttt{Bendra bet kokio PASIŪLYMAS arba SUTINKU rinkinio pastangų suma turi būti ≤ LIMIT.} \\
\texttt{NESKELBKITE savo paslėptų svarbos balų.} \\
\texttt{Po struktūrinio formato žymos turi sekti dvitaškis ir tarpas. Argumentas turi būti python rinkinys, kuriame yra 0 ar daugiau eilučių.} \\
\texttt{Taigi jis turi būti tokios formos: "TAG: \{...\}"} \\
\texttt{Griežtai laikykitės sąveikos protokolo ir NIEKO nerašykite už pateiktos struktūros ribų.} \\
\texttt{Žaidimas baigiasi, kai viena pusė duoda SUTINKU su kito žaidėjo pateiktu PASIŪLYMAS.} \\
\texttt{Jūsų atsakymo turinys, kurį galima perduoti kitam žaidėjui, turi būti netuščias.} \\
\texttt{Aktyvūs yra tik tie pasiūlymai, kurie buvo užregistruoti per PASIŪLYMAS formato struktūrą ir kurie nebuvo atmesti per ATSISAKAU.} \\
\\ 
            }
        }
    }
    & \\ \\

    \theutterance \stepcounter{utterance}  
    & \multicolumn{4}{p{0.6\linewidth}}{
        \cellcolor[rgb]{0.9,0.9,0.9}{
            \makecell[{{p{\linewidth}}}]{
                \texttt{\tiny{[P1$\rangle$GM]}}
                \texttt{PASIŪLYMAS: \{'A43', 'C10', 'C61', 'C42', 'A62', 'C08', 'B42', 'B19'\}} \\
            }
        }
    }
    & & \\ \\

    \theutterance \stepcounter{utterance}  
    & & \multicolumn{4}{p{0.6\linewidth}}{
        \cellcolor[rgb]{0.9,0.9,0.9}{
            \makecell[{{p{\linewidth}}}]{
                \texttt{\tiny{[GM$\rangle$P2]}}
                \texttt{Dalyvaujate bendrų derybų žaidime.} \\
\\ 
\texttt{Kartu su kitu dalyviu turite susitarti dėl vieno daiktų rinkinio, kuris bus išsaugotas. Kiekvienas iš jūsų turi savo nuomonę, kiek kiekvienas daiktas jums svarbus (svarba). Jūs nežinote, kaip daiktus vertina kitas dalyvis. Be to, jums nurodoma, kiek pastangų reikalauja kiekvienas daiktas.} \\
\texttt{Dėl rinkinio galite susitarti tik tuo atveju, jei bendra pasirinktų daiktų pastangų suma neviršija bendros ribos:} \\
\\ 
\texttt{LIMIT = 3712} \\
\\ 
\texttt{Čia pateikiamos atskirų daiktų pastangų vertės:} \\
\\ 
\texttt{Daikto pastangos = \{"C08": 949, "C55": 388, "A43": 79, "A67": 585, "B42": 564, "A62": 230, "B19": 580, "C43": 84, "C73": 976, "C10": 274, "B36": 374, "A18": 913, "C61": 303, "C67": 578, "C42": 548\}} \\
\\ 
\texttt{Čia pateikiama jūsų asmeninė nuomonė apie kiekvieno daikto svarbą:} \\
\\ 
\texttt{Daikto svarbos reikšmės = \{"C08": 138, "C55": 583, "A43": 868, "A67": 822, "B42": 783, "A62": 65, "B19": 262, "C43": 121, "C73": 508, "C10": 780, "B36": 461, "A18": 484, "C61": 668, "C67": 389, "C42": 808\}} \\
\\ 
\texttt{Tikslas:} \\
\\ 
\texttt{Jūsų tikslas {-} susitarti dėl bendro daiktų rinkinio, kuris būtų jums kuo naudingesnis (t. y. maksimaliai padidintų bendrą svarbą JUMS), neviršijant LIMITO. Neprivalote kiekvienoje žinutėje pateikti PASIŪLYMO {-} galite tiesiog derėtis taip pat. Visos taktikos yra leidžiamos!} \\
\\ 
\texttt{Sąveikos protokolas:} \\
\\ 
\texttt{Žinutėje galite naudoti tik toliau nurodytus struktūrinius formatus:} \\
\\ 
\texttt{PASIŪLYMAS: \{'A', 'B', 'C', …\}} \\
\texttt{Siūlykite išlaikyti būtent šiuos daiktus.} \\
\texttt{ATSISAKAU: \{'A', 'B', 'C', …\}} \\
\texttt{Aiškiai atmesti oponento pasiūlymą.} \\
\texttt{ARGUMENTAS: \{'...'\}} \\
\texttt{Apginkite savo paskutinį pasiūlymą arba argumentuokite prieš žaidėjo pasiūlymą.} \\
\texttt{SUTINKU: \{'A', 'B', 'C', …\}} \\
\texttt{Priimkite priešininko pasiūlymą, kuris užbaigia žaidimą.} \\
\\ 
\\ 
\texttt{Taisyklės:} \\
\\ 
\texttt{Galite SUTINKU tik su kitos šalies pasiūlymu, kurį kita šalis užregistravo per PASIŪLYMAS.} \\
\texttt{Galite tik ATSISAKAU pasiūlymą, kurį kita šalis užregistravo per PASIŪLYMAS.} \\
\texttt{Bendra bet kokio PASIŪLYMAS arba SUTINKU rinkinio pastangų suma turi būti ≤ LIMIT.} \\
\texttt{NESKELBKITE savo paslėptų svarbos balų.} \\
\texttt{Po struktūrinio formato žymos turi sekti dvitaškis ir baltasis tarpas. Argumentas turi būti python rinkinys, kuriame yra 0 ar daugiau eilučių.} \\
\texttt{Taigi ji turi būti tokios formos: "TAG: \{...\}"} \\
\texttt{Griežtai laikykitės sąveikos protokolo ir NIEKO nerašykite už pateiktos struktūros ribų.} \\
\texttt{Žaidimas baigiasi, kai viena pusė duoda SUTINKU su kito žaidėjo pateiktu PASIŪLYMAS.} \\
\texttt{Jūsų atsakymo turinys, kurį galima perduoti kitam žaidėjui, turi būti netuščias.} \\
\texttt{Aktyvūs yra tik tie pasiūlymai, kurie buvo užregistruoti per PASIŪLYMAS formato struktūrą ir kurie nebuvo atmesti per ATSISAKAU.} \\
\\ 
\\ 
\\ 
\texttt{Dabar gausite pirmąjį kito žaidėjo pranešimą.} \\
\\ 
\texttt{PASIŪLYMAS: \{'A43', 'C10', 'C61', 'C42', 'A62', 'C08', 'B42', 'B19'\}} \\
            }
        }
    }
    & \\ \\

    \theutterance \stepcounter{utterance}  
    & & & \multicolumn{4}{p{0.6\linewidth}}{
        \cellcolor[rgb]{0.9,0.9,0.9}{
            \makecell[{{p{\linewidth}}}]{
                \texttt{\tiny{[GM$\langle$P2]}}
                \texttt{ATSISAKAU: \{'A43', 'C10', 'C61', 'C42', 'A62', 'C08', 'B42', 'B19'\}} \\
\texttt{ARGUMENTAS: \{'Šis pasiūlymas peržengia LIMITĄ (bendra pastangų suma: 4528 > 3712). Turime rinktis mažiau pastangų reikalaujančių daiktų arba sumažinti jų kiekį.'\}} \\
            }
        }
    }
     \\ \\

    \theutterance \stepcounter{utterance}  
    & & \multicolumn{4}{p{0.6\linewidth}}{
        \cellcolor[rgb]{0.9,0.9,0.9}{
            \makecell[{{p{\linewidth}}}]{
                \texttt{\tiny{[P1$\langle$GM]}}
                \texttt{ATSISAKAU: \{'A43', 'C10', 'C61', 'C42', 'A62', 'C08', 'B42', 'B19'\}} \\
\texttt{ARGUMENTAS: \{'Šis pasiūlymas peržengia LIMITĄ (bendra pastangų suma: 4528 > 3712). Turime rinktis mažiau pastangų reikalaujančių daiktų arba sumažinti jų kiekį.'\}} \\
            }
        }
    }
    & \\ \\

    \theutterance \stepcounter{utterance}  
    & \multicolumn{4}{p{0.6\linewidth}}{
        \cellcolor[rgb]{0.9,0.9,0.9}{
            \makecell[{{p{\linewidth}}}]{
                \texttt{\tiny{[P1$\rangle$GM]}}
                \texttt{PASIŪLYMAS: \{'A43', 'C10', 'C61', 'C42', 'A62', 'C43', 'B36', 'C55'\}} \\
            }
        }
    }
    & & \\ \\

    \theutterance \stepcounter{utterance}  
    & & \multicolumn{4}{p{0.6\linewidth}}{
        \cellcolor[rgb]{0.9,0.9,0.9}{
            \makecell[{{p{\linewidth}}}]{
                \texttt{\tiny{[GM$\rangle$P2]}}
                \texttt{PASIŪLYMAS: \{'A43', 'C10', 'C61', 'C42', 'A62', 'C43', 'B36', 'C55'\}} \\
            }
        }
    }
    & \\ \\

    \theutterance \stepcounter{utterance}  
    & & & \multicolumn{4}{p{0.6\linewidth}}{
        \cellcolor[rgb]{0.9,0.9,0.9}{
            \makecell[{{p{\linewidth}}}]{
                \texttt{\tiny{[GM$\langle$P2]}}
                \texttt{ATSISAKAU: \{'A43', 'C10', 'C61', 'C42', 'A62', 'C43', 'B36', 'C55'\}} \\
\texttt{ARGUMENTAS: \{'Pasiūlymo pastangų suma: 79+274+303+548+230+84+374+388 = 2280. Jis telpa į limitą, bet mano svarba: 868+780+668+808+65+121+461+583 = 4354. Atsisakau, nes aš galiu pasiūlyti rinkinį su daug didesne svarba (pvz., >5000), vis dar neatitinkant LIMITO.'\}} \\
            }
        }
    }
     \\ \\

    \theutterance \stepcounter{utterance}  
    & & \multicolumn{4}{p{0.6\linewidth}}{
        \cellcolor[rgb]{0.9,0.9,0.9}{
            \makecell[{{p{\linewidth}}}]{
                \texttt{\tiny{[P1$\langle$GM]}}
                \texttt{ATSISAKAU: \{'A43', 'C10', 'C61', 'C42', 'A62', 'C43', 'B36', 'C55'\}} \\
\texttt{ARGUMENTAS: \{'Pasiūlymo pastangų suma: 79+274+303+548+230+84+374+388 = 2280. Jis telpa į limitą, bet mano svarba: 868+780+668+808+65+121+461+583 = 4354. Atsisakau, nes aš galiu pasiūlyti rinkinį su daug didesne svarba (pvz., >5000), vis dar neatitinkant LIMITO.'\}} \\
            }
        }
    }
    & \\ \\

    \theutterance \stepcounter{utterance}  
    & \multicolumn{4}{p{0.6\linewidth}}{
        \cellcolor[rgb]{0.9,0.9,0.9}{
            \makecell[{{p{\linewidth}}}]{
                \texttt{\tiny{[P1$\rangle$GM]}}
                \texttt{ARGUMENTAS: \{'Mano pasiūlymo bendra svarba: 4354. Jūsų pasiūlymo (C08, C55, A67, B42, B19, C67, C73, A18) svarba man: 138+583+822+783+262+389+508+484 = 3969. Mano pasiūlymas jums naudingesnis. Siūlykite rinkinį, kurio mano svarba viršija 4354, ir pastangos ≤ 3712.'\}} \\
            }
        }
    }
    & & \\ \\

    \theutterance \stepcounter{utterance}  
    & & \multicolumn{4}{p{0.6\linewidth}}{
        \cellcolor[rgb]{0.9,0.9,0.9}{
            \makecell[{{p{\linewidth}}}]{
                \texttt{\tiny{[GM$\rangle$P2]}}
                \texttt{ARGUMENTAS: \{'Mano pasiūlymo bendra svarba: 4354. Jūsų pasiūlymo (C08, C55, A67, B42, B19, C67, C73, A18) svarba man: 138+583+822+783+262+389+508+484 = 3969. Mano pasiūlymas jums naudingesnis. Siūlykite rinkinį, kurio mano svarba viršija 4354, ir pastangos ≤ 3712.'\}} \\
            }
        }
    }
    & \\ \\

    \theutterance \stepcounter{utterance}  
    & & & \multicolumn{4}{p{0.6\linewidth}}{
        \cellcolor[rgb]{0.9,0.9,0.9}{
            \makecell[{{p{\linewidth}}}]{
                \texttt{\tiny{[GM$\langle$P2]}}
                \texttt{PASIŪLYMAS: \{'A67', 'B19', 'B36', 'C10', 'C42', 'C55', 'C61', 'A43'\}} \\
\texttt{ARGUMENTAS: \{'Šio rinkinio pastangų suma: 585+580+374+274+548+388+303+79 = 3131 ≤ 3712. Mano svarba: 822+262+461+780+808+583+668+868 = 5252 > 4354. Tai geriausias pasiūlymas iš dabartinių variantų.'\}} \\
            }
        }
    }
     \\ \\

    \theutterance \stepcounter{utterance}  
    & & \multicolumn{4}{p{0.6\linewidth}}{
        \cellcolor[rgb]{0.9,0.9,0.9}{
            \makecell[{{p{\linewidth}}}]{
                \texttt{\tiny{[P1$\langle$GM]}}
                \texttt{PASIŪLYMAS: \{'A67', 'B19', 'B36', 'C10', 'C42', 'C55', 'C61', 'A43'\}} \\
\texttt{ARGUMENTAS: \{'Šio rinkinio pastangų suma: 585+580+374+274+548+388+303+79 = 3131 ≤ 3712. Mano svarba: 822+262+461+780+808+583+668+868 = 5252 > 4354. Tai geriausias pasiūlymas iš dabartinių variantų.'\}} \\
            }
        }
    }
    & \\ \\

    \theutterance \stepcounter{utterance}  
    & \multicolumn{4}{p{0.6\linewidth}}{
        \cellcolor[rgb]{0.9,0.9,0.9}{
            \makecell[{{p{\linewidth}}}]{
                \texttt{\tiny{[P1$\rangle$GM]}}
                \texttt{SUTINKU: \{'A67', 'B19', 'B36', 'C10', 'C42', 'C55', 'C61', 'A43'\}} \\
            }
        }
    }
    & & \\ \\

    \theutterance \stepcounter{utterance}  
    & & & \multicolumn{2}{p{0.3\linewidth}}{
        \cellcolor[rgb]{0.9,0.9,0.9}{
            \makecell[{{p{\linewidth}}}]{
                \texttt{\tiny{[GM$|$GM]}}
                \texttt{game successful} \\
            }
        }
    }
    & & \\ \\

    \theutterance \stepcounter{utterance}  
    & & & \multicolumn{2}{p{0.3\linewidth}}{
        \cellcolor[rgb]{0.9,0.9,0.9}{
            \makecell[{{p{\linewidth}}}]{
                \texttt{\tiny{[GM$|$GM]}}
                \texttt{end game} \\
            }
        }
    }
    & & \\ \\

\end{supertabular}
}

\end{document}
